\SourceLeaf{463}
\par\Needspace{8\baselineskip}\bigskip
\pdfbookmark[3]{Exposition, 14}{pallas-II-article14}
\begin{center}14.\end{center}
\noindent\hspace*{\parindent}Si, au lieu des deux forces \(mT\), \(mV\), nous en introduisons deux autres \(mR\), \(mS\), qui agissent également dans le plan de l’orbite, mais, la première dans la direction parallèle, la seconde dans la direction perpendiculaire à la ligne des apsides, ou plus exactement, que ces forces soient parallèles aux rayons de la sphère menés du centre vers les points \(P\), \(\Pi\), nous aurons
\[
\begin{aligned}
mR&=\chi\cos PQ=-\xi\cos XP-\eta\cos YP-\zeta\cos ZP\\
mS&=\chi\cos\Pi Q=-\xi\cos X\Pi-\eta\cos Y\Pi-\zeta\cos Z\Pi,
\end{aligned}
\]
où l’on peut mettre les valeurs des cosinus données par les équations 36—41. Nous aurons aussi
\[
\begin{aligned}
T&=R\cos(v-\tilde\omega)+S\sin(v-\tilde\omega)\\
V&=-R\sin(v-\tilde\omega)+S\cos(v-\tilde\omega).
\end{aligned}
\]
\par Nous ne donnerons ici que les résultats de la substitution de ces valeurs dans les différentes expressions des variations instantanées des élémens, en supprimant les détails des développemens, et en réduisant tout à la forme qui nous paraît la plus commode pour le calcul numérique:
\[
\begin{array}{rl}
72)\quad&\mathrm{d}a=-\dfrac{2a^3\sin(v-\tilde\omega)}{\cos\varphi}Rn\mathrm{d}t+\dfrac{2a^4\cos\varphi\mathbin{.}\cos E}{r}Sn\mathrm{d}t
\end{array}
\]
\[
\begin{array}{rl}
73)\quad&\mathrm{d}n=\dfrac{3aan\sin(v-\tilde\omega)}{\cos\varphi}Rn\mathrm{d}t-\dfrac{3a^3n\cos\varphi\mathbin{.}\cos E}{r}Sn\mathrm{d}t
\end{array}
\]
\[
\begin{array}{rl}
74)\quad&\begin{aligned}[t]\mathrm{d}\varphi&=-aa\cos E\mathbin{.}\sin(v-\tilde\omega)\mathbin{.}Rn\mathrm{d}t\\
&\quad+aa(1+\cos E\mathbin{.}\cos(v-\tilde\omega))Sn\mathrm{d}t\end{aligned}
\end{array}
\]
\[
\begin{array}{rl}
75)\quad&\begin{aligned}[t]\mathrm{d}\tilde\omega&=(1-\cos i)\mathrm{d}\AscendingNode\\
&\quad-aa\operatorname{cotang}\varphi\left(1+\dfrac{r\sin(v-\tilde\omega)^2}{p}\right)Rn\mathrm{d}t\\
&\quad+\dfrac{ar\sin2(v-\tilde\omega)}{\sin2\varphi}Sn\mathrm{d}t\\[4pt]
&=(1-\cos i)\mathrm{d}\AscendingNode-\dfrac{aa}{e}(\cos\varphi+\sin E\mathbin{.}\sin(v-\tilde\omega))Rn\mathrm{d}t\\
&\quad+\dfrac{aa}{e}\sin E\mathbin{.}\cos(v-\tilde\omega)\mathbin{.}Sn\mathrm{d}t\end{aligned}
\end{array}
\]
\[
\begin{array}{rl}
76)\quad&\begin{aligned}[t]\mathrm{d}\epsilon&=(1-\cos i)\mathrm{d}\AscendingNode\\
&\quad-\left\{\begin{aligned}2ar\cos(v-\tilde\omega)\\
{}+aa\operatorname{tang}\tfrac12\varphi(\cos\varphi+\sin E\mathbin{.}\sin(v-\tilde\omega))\end{aligned}\right\}Rn\mathrm{d}t\\
&\quad-aa\sin E(2\cos\varphi-\operatorname{tang}\tfrac12\varphi\mathbin{.}\cos(v-\tilde\omega))Sn\mathrm{d}t\\[4pt]
&=(1-\cos i)\mathrm{d}\AscendingNode\\
&\quad-\left\{\begin{aligned}2aa\cos E\\
{}-aa\operatorname{tang}\tfrac12\varphi(2+\cos\varphi\mathbin{.}\cos E\mathbin{.}\cos(v-\tilde\omega))\end{aligned}\right\}Rn\mathrm{d}t\\
&\quad-aa\sin E(2\cos\varphi-\operatorname{tang}\tfrac12\varphi\mathbin{.}\cos(v-\tilde\omega))Sn\mathrm{d}t.\end{aligned}
\end{array}
\]
